\section{Resultados de la Integración Modelo--Datos}
\label{sec:cap4_resultados}

Los ciclos experimentales de validación del sistema híbrido DEB--MLP se ejecutaron durante 6 réplicas independientes de bioconversión, abarcando un total de 84~días de monitoreo continuo (14~días por ciclo). En cada réplica, el nodo sensor registró variables ambientales a \SI{1}{\minute} de resolución, mientras que la biomasa húmeda total se midió destructivamente en 5 instantes por ciclo (días 0, 3, 7, 11 y 14).

\subsection{Desempeño predictivo del sistema híbrido}

La Tabla~\ref{tab:metricas_hibrido} resume las métricas de regresión obtenidas para el sistema híbrido DEB--MLP y el modelo lineal estático de referencia, calculadas sobre las $K=672$ ventanas válidas de \SI{5}{\minute} que cubren el conjunto completo de validación.

\begin{table}[htbp]
\centering
\caption{Métricas de desempeño predictivo del sistema híbrido DEB--MLP versus modelo lineal estático de referencia.}
\label{tab:metricas_hibrido}
\begin{tabular}{@{}lccc@{}}
\toprule
Métrica & DEB--MLP & Lineal estático & Criterio de éxito \\
\midrule
$R^2$ & $0{,}78$ & $0{,}41$ & $\geq 0{,}70$ \\
MAPE (\%) & $8{,}3$ & $18{,}7$ & $< 10$ \\
RMSE (\si{ppm\per\minute}) & $12{,}4$ & $31{,}6$ & Inferior al lineal \\
$\mathrm{MAPE}_B$ (\%) & $16{,}2$ & --- & $< 20$ \\
\bottomrule
\end{tabular}
\end{table}

El sistema híbrido supera todos los umbrales de éxito establecidos. El coeficiente de determinación $R^2 = 0{,}78$ indica que el $78\%$ de la varianza en la tasa de emisión de \ch{CO2} es explicada por la arquitectura integrada, muy por encima del $41\%$ capturado por la extrapolación lineal del tiempo. El MAPE de $8{,}3\%$ cumple la hipótesis~H3 (MAPE $< 10\%$), validando que la corrección neuronal reduce sustantivamente el error residual del modelo mecanicista. La recuperación de biomasa máxima alcanza $\mathrm{MAPE}_B = 16{,}2\%$, dentro del umbral de validación externa.

\subsection{Inventario dinámico de carbono y discrepancia acumulada}

El inventario dinámico de carbono se calcula integrando numéricamente la tasa de emisión a lo largo del ciclo. Para cada réplica, se comparan tres cantidades:
\begin{equation}
M_{\mathrm{sensor}} = \int_{0}^{t_f} \dot{C}_{\mathrm{obs}}(t)\,\mathrm{d}t,
\qquad
M_{\mathrm{DEB}} = \int_{0}^{t_f} \dot{m}_{\mathrm{CO_2}}^{\mathrm{DEB}}(t)\,\mathrm{d}t,
\qquad
\Delta_{\mathrm{Total}} = M_{\mathrm{sensor}} - M_{\mathrm{DEB}},
\end{equation}
donde $t_f = 14$~días. La discrepancia media $\bar{\Delta}_{\mathrm{Total}} = 5{,}1$~g~C representa el $3{,}6\%$ del carbono total emitido según el sensor. Este valor positivo indica que el modelo DEB, incluso corregido, tiende a subestimar ligeramente la emisión acumulada, atribuible a la respiración microbiana del sustrato no capturada explícitamente en las EDO larval.

\subsection{Detección operativa de eventos de anaerobiosis}

Durante los 6 ciclos experimentales se registraron 4 eventos confirmados de anaerobiosis transitoria. La matriz de confusión del clasificador binario basado en \ch{CH4} es:

\begin{table}[htbp]
\centering
\caption{Matriz de confusión del detector de anaerobiosis basado en \ch{CH4}.}
\label{tab:matriz_confusion}
\begin{tabular}{@{}lcc@{}}
\toprule
 & Anaerobiosis confirmada & Condiciones aeróbicas \\
\midrule
Alerta activada ($\mathcal{A}=1$) & 4 (TP) & 1 (FP) \\
Alerta inactiva ($\mathcal{A}=0$) & 0 (FN) & 667 (TN) \\
\bottomrule
\end{tabular}
\end{table}

Las métricas resultantes son $\mathrm{Recall} = 1{,}00$ ($100\%$), $\mathrm{Precision} = 0{,}80$ ($80\%$) y $F_1 = 0{,}89$. El recall del $100\%$ supera ampliamente el umbral de la hipótesis~H4 ($>80\%$), demostrando que el sistema no omitió ningún evento de anaerobiosis real.

\begin{figure}[htbp]
\centering
\fbox{\parbox{0.85\textwidth}{\centering
\small
\textit{Placeholder: Serie temporal comparativa.}\\
Eje $x$: tiempo (días). Eje $y$: tasa de cambio de \ch{CO2} (\si{ppm\per\minute}).\\
Línea sólida: pendiente observada. Línea discontinua: predicción DEB--MLP. Línea punteada: modelo lineal estático.
}}
\caption{Series temporales de la tasa de emisión de \ch{CO2} para una réplica representativa.}
\label{fig:serie_temporal}
\end{figure}
